%============================================================
% CHƯƠNG 8: ĐẠO CÔNG DÂN
%============================================================
\chapter{Đạo Công Dân (The Way of a Good Citizen)}

\begin{center}
  \textit{\color{truyenblue}``Ask not what your country can do for you --- ask what you can do for your country.''}\\
  \textit{(Đừng hỏi đất nước có thể làm gì cho bạn, hãy hỏi bạn có thể làm gì cho đất nước.)}\\[4pt]
  \small\color{truyengold}--- John F. Kennedy ---
\end{center}
\ngancach

\section{Chỗ Đỗ Xe}
\begin{truyen}{Chỗ Đỗ Xe}{Tuân Thủ}
\chuhoa{N}{ }gười đàn ông đỗ xe vào một chỗ hợp lệ dù còn chỗ thuận tiện hơn nhưng không được phép.\\
\textit{(The man parked in a legal spot even though a more convenient spot was available but not permitted.)}

Bạn anh trêu: ``Có ai nhìn đâu mà sợ.''\\
\textit{(His friend teased: ``Nobody is watching anyway.'')}

Anh cười: ``\textbf{Tao nhìn. Và tao không muốn là người mà chỉ làm đúng khi bị giám sát}.''\\
\textit{(He smiled: ``\textbf{I am watching. And I do not want to be someone who only does right when being watched}.'')}

Đây không phải là sợ bị phạt. Đây là \textbf{người ý thức rằng xã hội tốt hơn khi mỗi người tự nguyện giữ trật tự}.\\
\textit{(This was not fear of fines. It was \textbf{a person who understood that society improves when each individual voluntarily maintains order}.)}
\end{truyen}
\begin{baihoc}
  Đạo công dân không phải là tuân thủ luật pháp vì sợ bị phạt mà là vì hiểu rằng mỗi hành động của cá nhân đều có tác động đến toàn thể cộng đồng.
\end{baihoc}

\section{Cây Xanh Trước Nhà}
\begin{truyen}{Cây Xanh Trước Nhà}{Bảo Vệ Môi Trường}
\chuhoa{M}{ }ột hộ gia đình trồng cây xanh trước cửa và tưới đều đặn mỗi sáng.\\
\textit{(A household planted a tree in front of their home and watered it faithfully every morning.)}

Hàng xóm hỏi: ``Cây trồng trên đất công, người khác cũng được hưởng nhưng nhà anh phải tốn công chăm, có đáng không?''\\
\textit{(A neighbour asked: ``The tree is on public land, others benefit too but your household does all the work, is it worth it?'')}

``\textbf{Tôi không trồng cây cho nhà tôi. Tôi trồng cho con đường. Mà con đường là của tất cả mọi người, trong đó có tôi}.''\\
\textit{(``\textbf{I did not plant the tree for my house. I planted it for the street. And the street belongs to everyone, including me}.'')}

Năm năm sau, cả con đường mát rượi vì nhiều nhà học theo.\\
\textit{(Five years later, the whole street was shaded because many households followed his example.)}
\end{truyen}
\begin{baihoc}
  Công dân tốt không chờ người khác chăm sóc không gian chung. Họ chủ động đóng góp vì hiểu rằng không gian chung là không gian của chính họ.
\end{baihoc}

\section{Người Trả Lại Ví}
\begin{truyen}{Người Trả Lại Ví}{Trung Thực}
\chuhoa{M}{ }ột người lao động nhặt được chiếc ví dày trong công viên.\\
\textit{(A labourer found a fat wallet in the park.)}

Trong ví có số tiền bằng cả tháng lương của anh cộng với giấy tờ tùy thân của chủ nhân.\\
\textit{(Inside was money equal to a whole month's wages along with the owner's identification documents.)}

Anh đến đồn công an trình báo và nộp lại toàn bộ.\\
\textit{(He went to the police station to report it and handed everything over.)}

Chủ nhân tìm đến cảm ơn và đưa tiền thưởng. Anh từ chối:

``\textbf{Tôi trả lại vì đó là điều đúng, không phải để được thưởng. Nếu nhận thưởng cho việc đúng, lần sau tôi sẽ làm đúng vì tiền thay vì vì lương tâm}.''\\
\textit{(The owner came to thank him and offered a reward. He declined: ``\textbf{I returned it because it was right, not to be rewarded. If I accept payment for doing right, next time I will do right for money instead of for conscience}.'')}
\end{truyen}
\begin{baihoc}
  Đức thẳng thắn không cần phần thưởng bên ngoài. Lương tâm trong sạch là phần thưởng thiêng liêng nhất mà không ai có thể ban phát hay tước đoạt.
\end{baihoc}

\section{Hàng Xếp Thứ Tự}
\begin{truyen}{Hàng Xếp Thứ Tự}{Công Bằng}
\chuhoa{T}{ }rước cửa phòng khám, mọi người đều xếp hàng trật tự từ sáng sớm.\\
\textit{(Outside the clinic, everyone lined up orderly from early morning.)}

Một người đàn ông mặc vest đến trễ và định chen vào đầu hàng vì ``việc gấp''.\\
\textit{(A man in a suit arrived late and tried to cut to the front because of ``urgent business''.)}

Cụ già đứng đầu hàng nhẹ nhàng: ``Chúng tôi cũng có việc gấp. Đó là lý do chúng tôi đến sớm.''\\
\textit{(The elderly man at the front of the line said gently: ``We also have urgent matters. That is why we came early.'')}

Người mặc vest đỏ mặt và \textbf{đi về cuối hàng}.\\
\textit{(The suited man flushed and \textbf{walked to the back of the line}.)}

Một hành động nhỏ của cụ già đã \textbf{dạy cả hàng người về phẩm giá của sự công bằng}.\\
\textit{(The old man's small act \textbf{taught the entire line about the dignity of fairness}.)}
\end{truyen}
\begin{baihoc}
  Công dân tốt bảo vệ sự công bằng không chỉ cho mình mà cho toàn thể. Dũng cảm lên tiếng trước bất công nhỏ nhặt là hành động dân chủ bình thường nhất và quan trọng nhất.
\end{baihoc}

\section{Người Quét Rác Cuối Phố}
\begin{truyen}{Người Quét Rác Cuối Phố}{Trách Nhiệm}
\chuhoa{Ô}{ }ng lao công trên sáu mươi tuổi quét rác một con phố nhỏ suốt hai mươi năm.\\
\textit{(The more than sixty-year-old sanitation worker swept a small street for twenty years.)}

Người ta hỏi: ``Con phố này có quan trọng gì đâu ông.''\\
\textit{(People asked: ``This street is not even important, is it?'')}

Ông dừng chổi: ``\textbf{Đối với những người sống ở đây, nó là cả thế giới của họ}. Và mỗi sáng họ thức dậy thấy đường sạch, họ bắt đầu ngày mới tốt hơn. Đó là công việc của tôi.''\\
\textit{(He stopped his broom: ``\textbf{For the people who live here, it is their entire world}. And every morning they wake up to a clean street, they start their day better. That is my work.'')}

Khi ông về hưu, cả phố tập trung tiễn ông như tiễn một người thân.\\
\textit{(When he retired, the whole street gathered to see him off as if bidding farewell to a family member.)}
\end{truyen}
\begin{baihoc}
  Không có nghề nghiệp nào vô nghĩa nếu được thực hiện với ý thức trách nhiệm và tấm lòng vì cộng đồng. Giá trị của công việc không nằm ở danh tiếng mà nằm ở tác động tích cực đến cuộc sống người khác.
\end{baihoc}

\section{Lá Phiếu Của Người Già}
\begin{truyen}{Lá Phiếu Của Người Già}{Nghĩa Vụ Công Dân}
\chuhoa{C}{ }ụ bà chín mươi tuổi nhờ cháu chở đến điểm bầu cử dù sức đã rất yếu.\\
\textit{(The ninety-year-old grandmother asked her grandchild to drive her to the polling station despite being very weak.)}

``Bà ơi, bà mệt vậy không cần đi đâu con ơi.''\\
\textit{(``Grandmother, you are so tired, you do not have to go.'')}

Cụ lắc đầu: ``\textbf{Bà đã sống qua thời không có quyền bầu cử. Con không hiểu được quyền này quý giá như thế nào đâu. Bà còn đi được thì bà còn đi bỏ phiếu}.''\\
\textit{(She shook her head: ``\textbf{I lived through a time when there was no right to vote. You cannot understand how precious this right is. As long as I can walk, I will vote}.'')}

Cụ bỏ phiếu xong, ngồi nghỉ trước điểm bầu cử, mắt rưng rưng.\\
\textit{(After casting her ballot, she rested outside the polling station, eyes glistening.)}
\end{truyen}
\begin{baihoc}
  Những quyền công dân mà ta cho là hiển nhiên là kết quả của bao nhiêu đời người đã đấu tranh để có được. Trân trọng và sử dụng những quyền đó là cách tri ân xứng đáng nhất.
\end{baihoc}

\section{Đoạn Đường Không Tên}
\begin{truyen}{Đoạn Đường Không Tên}{Đóng Góp}
\chuhoa{M}{ }ột kỹ sư về hưu dùng tiền tiết kiệm cả đời tự làm con đường nối hai thôn nghèo.\\
\textit{(A retired engineer used his lifetime savings to personally build a road connecting two poor villages.)}

Ông làm một mình trong hai năm, không nhận sự hỗ trợ của chính quyền vì thủ tục quá phức tạp.\\
\textit{(He worked alone for two years, refusing government support because the paperwork was too complicated.)}

Khi xong, nhà chức trách hỏi ông muốn đặt tên gì cho con đường.\\
\textit{(When it was done, the authorities asked what he wanted to name the street.)}

``\textbf{Không cần tên. Đường không cần tên mới có người đi. Miễn là nó nối được hai làng lại với nhau là đủ rồi}.''\\
\textit{(``\textbf{No name needed. A road does not need a name to be walked on. As long as it connects the two villages, that is enough}.'')}
\end{truyen}
\begin{baihoc}
  Đóng góp cho cộng đồng không cần gắn tên mình vào đó. Công dân thực sự làm điều tốt vì điều tốt cần được làm, không phải vì cần người ghi nhận.
\end{baihoc}

\section{Người Cứu Đuối}
\begin{truyen}{Người Cứu Đuối}{Dũng Cảm}
\chuhoa{M}{ }ột thanh niên không biết bơi nhảy xuống sông cứu đứa trẻ đang chìm.\\
\textit{(A young man who could not swim jumped into the river to save a drowning child.)}

Cả hai được vớt lên bờ. Người thanh niên nằm thở dốc, gần bất tỉnh.\\
\textit{(Both were pulled to shore. The young man lay gasping, nearly unconscious.)}

Người xem hỏi: ``Anh không biết bơi sao anh dám nhảy?''\\
\textit{(Onlookers asked: ``You cannot swim, how did you dare jump?'')}

``\textbf{Tôi biết nếu không nhảy, đứa trẻ đó chết chắc. Tôi nhảy xuống thì còn có cơ hội. Đứa trẻ không quen biết gì với tôi, nhưng nó là người}.''\\
\textit{(``\textbf{I knew if I did not jump, the child would certainly die. If I jumped, there was still a chance. The child was a stranger to me, but the child was a person}.'')}
\end{truyen}
\begin{baihoc}
  Lòng dũng cảm thực sự không phải là không sợ hãi mà là hành động dù đang sợ hãi. Đạo làm người nằm ở khả năng nhìn người khác như chính bản thân mình.
\end{baihoc}

\section{Cuộc Gọi Báo Cháy}
\begin{truyen}{Cuộc Gọi Báo Cháy}{Trách Nhiệm Xã Hội}
\chuhoa{L}{ }úc ba giờ sáng, người phụ nữ nhìn thấy khói bay ra từ nhà hàng xóm.\\
\textit{(At three in the morning, the woman noticed smoke coming from her neighbour's house.)}

Bà nghĩ mãi xem có nên gọi điện không vì sợ giật mình người ta.\\
\textit{(She wondered for a moment whether to call, afraid of startling the neighbours.)}

Rồi bà gọi ngay lập tức.\\
\textit{(Then she called immediately.)}

Cả nhà hàng xóm chạy thoát ra kịp trước khi lửa bùng lớn.\\
\textit{(The entire household escaped before the fire spread.)}

Chủ nhà cảm ơn: ``Nếu không có bà, chúng tôi không kịp chạy.'' Bà trả lời: ``\textbf{Tôi không thể ngủ trong khi hàng xóm đang nguy hiểm}.''\\
\textit{(The house owner thanked her: ``Without you, we would not have escaped.'' She replied: ``\textbf{I cannot sleep while my neighbours are in danger}.'')}
\end{truyen}
\begin{baihoc}
  Công dân tốt không thờ ơ trước nguy hiểm của người xung quanh. Cộng đồng an toàn và bình yên được xây nên từ hàng triệu khoảnh khắc nhỏ khi mỗi người chọn quan tâm thay vì quay mặt đi.
\end{baihoc}

\section{Tờ Đơn Kiến Nghị}
\begin{truyen}{Tờ Đơn Kiến Nghị}{Lên Tiếng}
\chuhoa{N}{ }gôi trường tiểu học gần nhà xuống cấp nguy hiểm nhưng không ai dám lên tiếng vì sợ bị phiền phức.\\
\textit{(The nearby primary school was dangerously run-down but no one dared speak up for fear of complications.)}

Một phụ huynh gõ cửa từng nhà trong khu, vận động ký tên vào tờ đơn kiến nghị.\\
\textit{(One parent knocked on every door in the neighbourhood, mobilising signatures for a petition.)}

Nhiều người từ chối. Bà không bỏ cuộc.\\
\textit{(Many refused. She did not give up.)}

Ba tháng sau, trường được sửa chữa.\\
\textit{(Three months later, the school was repaired.)}

``Sao bà không sợ bị để ý?'' – ``\textbf{Tôi sợ. Nhưng tôi sợ gửi con đến một ngôi trường có thể sập xuống hơn}.''\\
\textit{(``Why were you not afraid of drawing attention?'' – ``\textbf{I was afraid. But I was more afraid of sending my child to a building that could collapse}.'')}
\end{truyen}
\begin{baihoc}
  Công dân có trách nhiệm không chỉ im lặng chịu đựng mà còn có quyền và nghĩa vụ lên tiếng đòi hỏi những điều đúng đắn cho cộng đồng mình sống trong đó.
\end{baihoc}
